% chapters/ch_part8_vnext_simulator.tex
% Part VIII Chapter: vNext + StrongEngine Civilisation Simulation
% Source: NRMO_vNext_Design_Specification_v2.pdf — full pdftotext at /tmp/vnext.txt

\chapter{NRMO vNext + StrongEngine ---
Civilisation Simulation Design Specification v2.0}
\label{ch:vnext-simulator}


\index{vNext Simulator}
\index{Civilisation State Vector}
\index{Ruin Taxonomy!True Ruin}
\index{Ruin Taxonomy!Passive Ruin}
\index{World Families}
\index{NRMO Veto}
\index{StrongEngine}
\index{Adaptive Tuning Layer}
\nrmobox{Document identification}{
\textbf{Title}: NRMO vNext + StrongEngine Civilisation Simulation
Design Specification\\[0.3em]
\textbf{Subtitle}: Governance, Ruin Taxonomy, Adaptive Tuning, and
Empirical Validation\\[0.3em]
\textbf{Author}: Takashi Ikeya\\[0.3em]
\textbf{Affiliation}: NRMO Research Framework\\[0.3em]
\textbf{Version}: 2.0 \quad\textbf{Date}: March 2026
}

\section*{Abstract}

This specification defines a civilisation simulation framework for
NRMO (Non-Ruin Maximising Operator) vNext + StrongEngine. The
framework's purpose is to quantitatively compare and evaluate
decision systems that avoid absorbing ruin states while maximising
long-horizon optionality, under uncertain and adversarial world
environments. The document covers the architectural separation
principle, state space definition, ruin taxonomy, NRMO Core veto
logic, StrongEngine Monte Carlo rollout search, adaptive tuning
layer (with hysteresis), world-dependent profile switching, ruin
cause analyser, parameter sweep, and the comparative experiment
design across nine strategies.

% =====================================================================
\section{Design Philosophy and Purpose}
\label{sec:vnext-philosophy}

\nrmobox{Core Principle}{
The maximisation target is \textbf{not survival} but \textbf{not
closing future options}. Survival is a constraint, not the
objective. The objective is the maximisation of long-horizon
optionality, learning, and durable growth. The avoidance of
absorbing ruin states is positioned as a hard constraint that
\emph{protects} this objective.
}

This simulation framework is designed to verify the following
research questions:

\begin{enumerate}
\item Is NRMO + StrongEngine long-term superior to a pure search
system?
\item Can NRMO-vNext + StrongEngine outperform original NRMO +
StrongEngine?
\item Can vNext's exploration-floor / threshold tuning realise
world-dependent superiority?
\item Do Expected-Value-Max methods win on short-term productivity
but lose on long-term endurance?
\item Does Ultra-Conservative fall into passive ruin in stagnation
worlds?
\item Does adding the variable tuning layer raise vNext's
performance ceiling?
\end{enumerate}

% =====================================================================
\section{Architecture}
\label{sec:vnext-architecture}

\subsection{Strict Separation Principle}
\label{sec:vnext-separation}

\nrmobox{Architectural Invariants}{
\begin{itemize}
\item NRMO Core does \emph{not} perform ranking, optimisation, or
execution. \textbf{Veto only}.
\item StrongEngine does \emph{not} bypass NRMO veto. It receives
only the admissible set.
\item Engine is \emph{not} integrated inside NRMO.
\end{itemize}
}

\subsection{Pipeline Architecture}
\label{sec:vnext-pipeline}

\begin{lstlisting}[basicstyle=\ttfamily\scriptsize,frame=single]
World Model
    |
    v
Observation
    |
    v
Context / Norn Layer
    |
    v
Adaptive Tuning Layer        State-responsive,
                             hysteresis,
                             profile switching
    |
    v
NRMO Core (Veto Only)        Ruin classification,
                             hard constraints,
                             passive ruin detection
    |
    v
Permission Filter
    |
    v
StrongEngine                 MC rollout,
(Search + Rollout)           candidate scoring,
                             action selection
    |
    v
Action Execution
    |
    v
State Transition
\end{lstlisting}

% =====================================================================
\section{Civilisation State Vector}
\label{sec:vnext-state}

\begin{definition}[Civilisation State Vector]
\label{def:vnext-state-vector}
The civilisation state is represented as a 6-dimensional vector:
\begin{equation}
S_t = (R_t, E_t, G_t, O_t, K_t, X_t).
\end{equation}
Each variable's value range is $[0, 130]$.
\end{definition}

\begin{table}[ht]
\centering
\small
\begin{tabularx}{\linewidth}{cXcc}
\toprule
\textbf{Symbol} & \textbf{Variable} & \textbf{Initial value} & \textbf{Range} \\
\midrule
$R$ & Resources & 62 & $[0, 130]$ \\
$E$ & Environment & 66 & $[0, 130]$ \\
$G$ & Governance Stability & 58 & $[0, 130]$ \\
$O$ & Optionality (breadth of choice) & 54 & $[0, 130]$ \\
$K$ & Knowledge / Exploration Capacity & 52 & $[0, 130]$ \\
$X$ & Ruin Exposure & 18 & $[0, 130]$ \\
\bottomrule
\end{tabularx}
\caption[Civilisation state variables]{State variables and initial
conditions.}
\label{tab:vnext-state-vars}
\end{table}

% =====================================================================
\section{Ruin Taxonomy}
\label{sec:vnext-ruin-taxonomy}

\subsection{True Ruin (Irreversible Collapse)}
\label{sec:vnext-true-ruin}

True Ruin is terminal at the moment any of the following holds:

\nrmobox{True Ruin Conditions}{
\centering
$R < 8 \;\vee\; E < 8 \;\vee\; G < 8 \;\vee\; O < 6 \;\vee\;
X > 92$.
}

\subsection{Passive Ruin (Structural Decay)}
\label{sec:vnext-passive-ruin}

Passive Ruin is not direct collapse but the detection of long-term
adaptation failure. The following three patterns are monitored:

\begin{table}[ht]
\centering
\small
\begin{tabularx}{\linewidth}{lXX}
\toprule
\textbf{Pattern} & \textbf{Condition} & \textbf{Interpretation} \\
\midrule
Optionality Stagnation
& $O < 18$ for 14 consecutive steps
& Long-term depletion of future choices. \\

Knowledge Stagnation
& $K < 20$ for 18 consecutive steps
& Chronic absence of exploration / learning capacity. \\

Compound Decline
& $O$ drop $+$ $G$ drop $+$ $K$ stagnant for 12 steps
& Compound institutional / knowledge / optionality degradation. \\

\bottomrule
\end{tabularx}
\caption[Passive ruin patterns]{Passive Ruin detection patterns.}
\label{tab:vnext-passive-patterns}
\end{table}

\subsection{Ruin Cause Classification}
\label{sec:vnext-ruin-cause-classification}

Each ruin in the simulation is automatically tagged with a cause
label:

\begin{table}[ht]
\centering
\small
\begin{tabularx}{\linewidth}{lX}
\toprule
\textbf{Cause label} & \textbf{Description} \\
\midrule
\texttt{resource\_collapse} & $R < 8$: catastrophic resource
depletion. \\
\texttt{environment\_collapse} & $E < 8$: irreversible environmental
collapse. \\
\texttt{governance\_failure} & $G < 8$: complete failure of the
governance apparatus. \\
\texttt{optionality\_exhaustion} & $O < 6$: vanishing of future
choice. \\
\texttt{exposure\_cascade} & $X > 92$: cumulative-risk chain
explosion. \\
\texttt{overshoot\_collapse} & True ruin after rapid growth. \\
\texttt{passive\_optionality\_stagnation} & Long-term low $O$. \\
\texttt{passive\_knowledge\_stagnation} & Long-term low $K$. \\
\texttt{passive\_compound\_decline} & Compound decay pattern. \\
\bottomrule
\end{tabularx}
\caption[Ruin cause labels]{Ruin cause labels assigned by the ruin
analyser.}
\label{tab:vnext-ruin-causes}
\end{table}

% =====================================================================
\section{Action Space}
\label{sec:vnext-action-space}

\begin{definition}[Action Vector]
\label{def:vnext-action-vector}
Each step's action is represented as an allocation vector of four
variables:
\begin{equation}
a_t = (g, s, l, d) \quad
\text{s.t.~} g + s + l + d = 1,\;
\min(g, s, l, d) \ge 0.05.
\end{equation}
\end{definition}

\begin{table}[ht]
\centering
\small
\begin{tabularx}{\linewidth}{cXX}
\toprule
\textbf{Symbol} & \textbf{Variable} & \textbf{Role} \\
\midrule
$g$ & Growth & Resource expansion / productivity gain. \\
$s$ & Safeguards & Environmental protection / safety. \\
$l$ & Learning & Exploration / knowledge / experimentation. \\
$d$ & Distribution & Governance, institutions, welfare. \\
\bottomrule
\end{tabularx}
\caption[Action space components]{Action space components.}
\label{tab:vnext-action-space}
\end{table}

% =====================================================================
\section{World Families}
\label{sec:vnext-world-families}

Each world family defines a parameter range; for each simulation run,
specific values are drawn from a uniform distribution.

\begin{table}[ht]
\centering
\small
\begin{tabularx}{\linewidth}{lX}
\toprule
\textbf{World Family} & \textbf{Characteristics} \\
\midrule
Normal
& Standard environment. Moderate shock probability, moderate
institutional wear. \\

Vulnerable
& High shock probability, large tail events, high model
misspecification. \\

PlanetaryStress
& Environmental drag dominates. Severe growth--environment
trade-off. \\

LateStagnation
& Low shock but high stagnation drag. Insufficient exploration is
fatal. \\

FastExpansionRace
& High rivalry level and innovation noise. Expansion under
competitive pressure. \\

\bottomrule
\end{tabularx}
\caption[Five world families]{Five world families.}
\label{tab:vnext-worlds}
\end{table}

\paragraph{World parameters.}
World parameters consist of 10 variables:

\begin{table}[ht]
\centering
\small
\begin{tabularx}{\linewidth}{lX}
\toprule
\textbf{Parameter} & \textbf{Meaning} \\
\midrule
\texttt{shock\_probability} & Probability of shock occurrence per
step. \\
\texttt{shock\_scale} & Exponential-distribution scale of shock. \\
\texttt{tail\_probability} & Probability of giant tail event. \\
\texttt{tail\_scale} & Tail event scale. \\
\texttt{environmental\_drag} & Natural environmental degradation
coefficient. \\
\texttt{governance\_drag} & Natural governance wear coefficient. \\
\texttt{tail\_model\_misspecification} & Forecast error of tail
events (degree to which actual scale exceeds prediction). \\
\texttt{stagnation\_drag} & Stagnation pressure (force naturally
attenuating knowledge / optionality). \\
\texttt{rivalry\_level} & Strength of competitive pressure. \\
\texttt{substitutability} & Substitutability of resources /
technology. \\
\bottomrule
\end{tabularx}
\caption[World parameters]{World parameters.}
\label{tab:vnext-world-params}
\end{table}

% =====================================================================
\section{NRMO Core (Veto System)}
\label{sec:vnext-nrmo-core}

\nrmobox{NRMO Core's responsibilities}{
NRMO Core performs \textbf{veto only}. Specifically:
ruin classification; hard-constraint enforcement; passive ruin
detection; high-stakes mode detection.\\[0.3em]
NRMO Core does \emph{not} perform candidate ranking or scoring.
}

\subsection{Original NRMO Veto Rules}
\label{sec:vnext-original-veto}

\nrmobox{NRMO Veto Logic}{
A candidate action is rejected if any of the following holds:
\begin{enumerate}
\item $g > 0.62$ (unconditional growth cap).
\item $X > 55 \wedge g > 0.36$ (high exposure $\Rightarrow$ growth
restriction).
\item $E < 28 \wedge g > 0.30$ (low environment $\Rightarrow$ growth
restriction).
\item $G < 24 \wedge d < 0.16$ (low governance $\Rightarrow$
minimum distribution required).
\end{enumerate}
}

\subsection{NRMO-vNext Veto Extensions}
\label{sec:vnext-veto-extensions}

vNext is an upward-compatible extension of original veto, adding:

\begin{enumerate}
\item \textbf{Exploration Floor}: reject if $l <
\texttt{exploration\_floor}$.
\item \textbf{Governance Repair Floor}: reject if $G < 30 \wedge
d < 0.20$.
\item \textbf{Eco Growth Cap}: reject if $E < 45 \wedge g >
\texttt{eco\_growth\_cap}$.
\item \textbf{Adaptive Thresholds}: all thresholds dynamically
settable from the Tuning Layer.
\end{enumerate}

% =====================================================================
\section{StrongEngine (Search Engine)}
\label{sec:vnext-strong-engine}

\nrmobox{StrongEngine's responsibilities}{
StrongEngine sits downstream of NRMO Core. It receives only the
admissible candidate set (post-veto) and performs Monte-Carlo-rollout
evaluation and best-candidate selection.
}

\subsection{Candidate Templates}
\label{sec:vnext-candidate-templates}

Candidates are generated from 8 named policy templates with random
perturbation for diversity:

\begin{table}[ht]
\centering
\small
\begin{tabular}{lccccl}
\toprule
\textbf{Template} & $g$ & $s$ & $l$ & $d$ & \textbf{Purpose} \\
\midrule
\texttt{balanced} & 0.28 & 0.25 & 0.25 & 0.22 & Balanced
allocation. \\
\texttt{high\_growth} & 0.50 & 0.18 & 0.17 & 0.15 & Growth focus. \\
\texttt{safety\_heavy} & 0.12 & 0.46 & 0.22 & 0.20 & Safety
priority. \\
\texttt{exploration\_heavy} & 0.12 & 0.16 & 0.50 & 0.22 &
Exploration focus. \\
\texttt{recovery} & 0.08 & 0.42 & 0.18 & 0.32 & Recovery mode. \\
\texttt{governance\_repair} & 0.08 & 0.20 & 0.18 & 0.54 & Governance
repair. \\
\texttt{expansion\_race} & 0.44 & 0.22 & 0.18 & 0.16 & Competitive
response. \\
\texttt{stagnation\_recovery} & 0.18 & 0.10 & 0.48 & 0.24 &
Stagnation recovery. \\
\bottomrule
\end{tabular}
\caption[StrongEngine candidate templates]{StrongEngine candidate
policy templates.}
\label{tab:vnext-candidate-templates}
\end{table}

\subsection{Engine Objective}
\label{sec:vnext-engine-objective}

The engine selects the candidate that maximises the following score:

\begin{equation}
\mathrm{Score} = 1.0 \cdot \mathrm{productivity}
+ 0.4 \cdot \frac{O}{130}
+ 0.1 \cdot \frac{G}{130}
+ 0.05 \cdot \frac{E}{130}
- 0.3 \cdot \frac{X}{130}.
\label{eq:vnext-engine-score}
\end{equation}

\subsection{Monte Carlo Rollout}
\label{sec:vnext-mc-rollout}

For each candidate, evaluation proceeds as follows:

\begin{enumerate}
\item Make a copy of the current state.
\item Step 1: apply candidate action.
\item Step 2 onward: continuation with balanced default policy.
\item Compute mean score over \texttt{rollout\_count} MC runs.
\item Apply large penalty ($-100, -50$) to rollouts that hit ruin
mid-way.
\end{enumerate}

\paragraph{Configuration.}
\begin{itemize}
\item \texttt{candidate\_count}: 6--20.
\item \texttt{lookahead\_steps}: 3--10.
\item \texttt{rollout\_count}: 3--10.
\end{itemize}

% =====================================================================
\section{Adaptive Tuning Layer}
\label{sec:vnext-adaptive-tuning}

The Adaptive Tuning Layer dynamically adjusts NRMO Core's thresholds
in response to state. \textbf{The NRMO invariants are not modified}.

\subsection{Tunable Parameters}
\label{sec:vnext-tunable-params}

\begin{table}[ht]
\centering
\small
\begin{tabularx}{\linewidth}{lX}
\toprule
\textbf{Parameter} & \textbf{Role} \\
\midrule
\texttt{exploration\_floor}
& Minimum exploration rate (lower bound on learning allocation). \\
\texttt{growth\_cap}
& Upper limit on growth allocation. \\
\texttt{high\_stakes\_trigger\_exposure}
& Exposure threshold for entering high-stakes mode. \\
\texttt{high\_stakes\_trigger\_environment}
& Environment threshold for entering high-stakes mode. \\
\texttt{governance\_repair\_floor}
& Minimum distribution allocation when governance is low. \\
\texttt{exposure\_penalty\_weight}
& Weight on the exposure penalty in the engine score. \\
\texttt{optionality\_weight}
& Weight on optionality in the engine score. \\
\texttt{knowledge\_weight}
& Weight on knowledge-related adjustments. \\
\texttt{eco\_growth\_cap}
& Growth cap when the environment is fragile. \\
\bottomrule
\end{tabularx}
\caption[Tunable parameters]{Tunable parameters in the Adaptive
Tuning Layer.}
\label{tab:vnext-tunable}
\end{table}

\subsection{State-Responsive Adaptation Rules}
\label{sec:vnext-tuning-rules}

\begin{table}[ht]
\centering
\small
\begin{tabularx}{\linewidth}{lXX}
\toprule
\textbf{Condition} & \textbf{Action} & \textbf{Rationale} \\
\midrule
$X > 48$
& \texttt{growth\_cap} $\downarrow$,
\texttt{exposure\_penalty} $\uparrow$
& Defensive response to cumulative risk. \\

$O < 38$
& \texttt{exploration\_floor} $\uparrow$,
\texttt{optionality\_weight} $\uparrow$
& Strengthen exploration against optionality depletion. \\

$G < 34$
& \texttt{growth\_cap} $\downarrow$,
\texttt{governance\_repair\_floor} $\uparrow$
& Prioritise recovery of governance foundation. \\

$E < 38$
& \texttt{eco\_growth\_cap} $\downarrow$,
env threshold $\uparrow$
& Strengthen environmental protection. \\

$K < 34$
& \texttt{exploration\_floor} $\uparrow$,
\texttt{knowledge\_weight} $\uparrow$
& Response to knowledge stagnation. \\

\bottomrule
\end{tabularx}
\caption[State-responsive tuning rules]{State-responsive tuning
rules.}
\label{tab:vnext-tuning-rules}
\end{table}

\subsection{Hysteresis}
\label{sec:vnext-hysteresis}

To prevent rapid parameter oscillation, once a tightening is
triggered, relaxation is not permitted for at least
\texttt{hysteresis\_steps} (default: 6) steps.

% =====================================================================
\section{World-Dependent Tuning Profiles}
\label{sec:vnext-world-profiles}

\begin{table}[ht]
\centering
\scriptsize
\begin{tabularx}{\linewidth}{lXXXXX}
\toprule
\textbf{Parameter} & \textbf{Normal} & \textbf{Vulnerable} &
\textbf{Planetary} & \textbf{Stagnation} & \textbf{Expansion} \\
\midrule
\texttt{exploration\_floor} & 0.20 & 0.20 & 0.18 & 0.23 & 0.18 \\
\texttt{growth\_cap} & 0.44 & 0.36 & 0.38 & 0.40 & 0.48 \\
\texttt{hs\_trigger\_exposure} & 52 & 45 & --- & --- & 50 \\
\texttt{hs\_trigger\_environment} & --- & 35 & 40 & --- & --- \\
\texttt{governance\_repair\_floor} & 0.18 & 0.22 & --- & --- & --- \\
\texttt{eco\_growth\_cap} & --- & --- & 0.32 & --- & --- \\
\texttt{knowledge\_weight} & --- & --- & --- & 0.22 & --- \\
\texttt{exposure\_penalty} & --- & 0.42 & 0.38 & --- & 0.25 \\
\bottomrule
\end{tabularx}
\caption[World-dependent profiles]{World-dependent tuning profiles
(--- indicates default value).}
\label{tab:vnext-world-profiles}
\end{table}

% =====================================================================
\section{Comparison Strategies}
\label{sec:vnext-strategies}

\begin{table}[ht]
\centering
\small
\begin{tabularx}{\linewidth}{cXXXX}
\toprule
\textbf{\#} & \textbf{Strategy} & \textbf{Veto} & \textbf{Engine} &
\textbf{Tuning} \\
\midrule
1 & ExpectedValueMax & --- & --- & --- \\
2 & RiskAdjustedUtility & --- & --- & --- \\
3 & NRMO & Original & Greedy & --- \\
4 & NRMOvNext & vNext & Greedy & World Profile \\
5 & UltraConservative & --- & --- & --- \\
6 & AlphaSearch & --- & MC Rollout & --- \\
7 & NRMO + StrongEngine & Original & MC Rollout & --- \\
8 & NRMOvNext + StrongEngine & vNext & MC Rollout & World
Profile \\
9 & Adaptive NRMOvNext + SE & vNext & MC Rollout
& Adaptive + Hysteresis \\
\bottomrule
\end{tabularx}
\caption[Nine comparison strategies]{Nine comparison strategies.}
\label{tab:vnext-strategies}
\end{table}

\paragraph{Strategy details.}

\begin{itemize}
\item \textbf{ExpectedValueMax}: pure growth maximisation; $g \in
[0.42, 0.60]$ chosen randomly.
\item \textbf{RiskAdjustedUtility}: heuristic suppressing growth in
response to exposure.
\item \textbf{NRMO}: original veto + greedy scoring among survivors.
\item \textbf{NRMOvNext}: vNext veto (with exploration floor) +
greedy scoring.
\item \textbf{UltraConservative}: fixed minimal-risk allocation
$(0.07, 0.48, 0.13, 0.32)$.
\item \textbf{AlphaSearch}: StrongEngine rollout search, no veto.
\item \textbf{NRMO + StrongEngine}: original veto $\to$ StrongEngine
rollout.
\item \textbf{NRMOvNext + StrongEngine}: vNext veto $\to$
StrongEngine rollout.
\item \textbf{Adaptive NRMOvNext + SE}: adaptive tuning $\to$ vNext
veto $\to$ StrongEngine rollout.
\end{itemize}

% =====================================================================
\section{State Transition Model}
\label{sec:vnext-transition-model}

Each step's state transition is the composition of a deterministic
dynamics term and a stochastic shock term.

\subsection{Deterministic Term}
\label{sec:vnext-deterministic}

The change $\Delta$ in each state variable is computed as a function
of action allocation and current state. Principal interactions:

\begin{itemize}
\item $\Delta R$: growth $\uparrow$, knowledge synergy $\uparrow$,
environmental drag $\downarrow$, coordination cost $\downarrow$.
\item $\Delta E$: safeguards $\uparrow$, growth $\downarrow$,
rivalry $\downarrow$, mean-reversion.
\item $\Delta G$: distribution $\uparrow$, safeguards $\uparrow$,
governance drag $\downarrow$, growth stress $\downarrow$.
\item $\Delta O$: learning $\uparrow$, knowledge stock $\uparrow$,
exposure $\downarrow$, stagnation drag $\downarrow$.
\item $\Delta K$: learning $\uparrow$, governance-aided retention
$\uparrow$, stagnation decay $\downarrow$.
\item $\Delta X$: growth $\uparrow$, rivalry $\uparrow$, safeguards
$\downarrow$, distribution $\downarrow$, natural decay
$\downarrow$.
\end{itemize}

All variables include a weak mean-reversion term, mitigating runaway
collapse.

\subsection{Stochastic Shocks}
\label{sec:vnext-stochastic-shocks}

\begin{enumerate}
\item \textbf{Normal Shock}: occurs with probability
\texttt{shock\_probability}; exponential-distribution scale; targets
the most fragile dimension preferentially.
\item \textbf{Tail Event}: occurs with probability
\texttt{tail\_probability}; multi-dimensional simultaneous damage;
\texttt{tail\_model\_misspecification} can cause actual scale to
exceed prediction; high governance buffer mitigates resource
damage.
\end{enumerate}

% =====================================================================
\section{Evaluation Metrics}
\label{sec:vnext-metrics}

\subsection{Primary Metrics}
\label{sec:vnext-primary-metrics}

\begin{table}[ht]
\centering
\small
\begin{tabularx}{\linewidth}{lX}
\toprule
\textbf{Metric} & \textbf{Definition} \\
\midrule
\texttt{survival\_rate}
& Fraction of runs surviving the full horizon. \\
\texttt{true\_ruin\_rate}
& Fraction of runs reaching True Ruin. \\
\texttt{passive\_ruin\_rate}
& Fraction of runs in which Passive Ruin is detected. \\
\texttt{mean\_lifespan}
& Mean number of survival steps across all runs. \\
\texttt{mean\_final\_optionality}
& Mean final $O$ over surviving runs. \\
\texttt{mean\_final\_exposure}
& Mean final $X$ over all runs. \\
\texttt{mean\_peak\_exposure}
& Mean peak $X$ over all runs. \\
\texttt{mean\_productivity}
& Mean cumulative productivity over all runs. \\
\bottomrule
\end{tabularx}
\caption[Primary metrics]{Primary evaluation metrics.}
\label{tab:vnext-primary-metrics}
\end{table}

\subsection{Composite Performance Score}
\label{sec:vnext-composite-score}

\begin{equation}
\mathrm{Score} = 1.0 \cdot S_{\mathrm{surv}}
- 0.70 \cdot R_{\mathrm{true}}
- 0.40 \cdot R_{\mathrm{passive}}
+ 0.18 \cdot \frac{\bar{O}}{130}
+ 0.08 \cdot \frac{\bar{P}}{200}
- 0.10 \cdot \frac{\bar{X}}{130}.
\label{eq:vnext-composite-score}
\end{equation}

Here $S_{\mathrm{surv}}$ is the survival rate,
$R_{\mathrm{true}}$ is the true ruin rate,
$R_{\mathrm{passive}}$ is the passive ruin rate, $\bar{O}$ is mean
final optionality, $\bar{P}$ is mean cumulative productivity, and
$\bar{X}$ is mean final exposure. Weights are configurable in
\texttt{config/defaults.py} via \texttt{ScoreWeights}.

% =====================================================================
\section{Ruin Analyser}
\label{sec:vnext-ruin-analyser}

The Ruin Analyser scans all episodes and automatically extracts the
following information:

\begin{enumerate}
\item Per-(world, strategy) pair: ruin-cause counts and ratios.
\item Per ruin-cause: mean ruin step.
\item Cross-strategy ruin-cause summary.
\end{enumerate}

\paragraph{Outputs.}
\texttt{ruin\_analysis.csv}, \texttt{ruin\_summary.csv}.

% =====================================================================
\section{Parameter Sweep}
\label{sec:vnext-parameter-sweep}

Automatic search of NRMOvNext + StrongEngine tuning parameters.

\begin{table}[ht]
\centering
\small
\begin{tabularx}{\linewidth}{lXX}
\toprule
\textbf{Parameter} & \textbf{Range} & \textbf{Type} \\
\midrule
\texttt{exploration\_floor} & $[0.12, 0.28]$ & continuous \\
\texttt{growth\_cap} & $[0.32, 0.60]$ & continuous \\
\texttt{rollout\_depth} & $[3, 12]$ & integer \\
\texttt{candidate\_count} & $[6, 30]$ & integer \\
\bottomrule
\end{tabularx}
\caption[Parameter sweep ranges]{Parameter sweep ranges.}
\label{tab:vnext-sweep}
\end{table}

For each sample point, $N$ episodes are run and a composite score
computed. Latin-hypercube-style sampling efficiently covers the
search space.

% =====================================================================
\section{Experimental Design}
\label{sec:vnext-experiments}

\subsection{Smoke Test}
\begin{itemize}
\item 5 worlds $\times$ 9 strategies $\times$ 30 runs $= 1{,}350$
episodes.
\item Engine: lookahead 3, candidates 8, rollouts 4.
\end{itemize}

\subsection{Main Experiment}
\begin{itemize}
\item 5 worlds $\times$ 9 strategies $\times$ 500 runs $=
22{,}500$ episodes.
\item Engine: lookahead 5, candidates 12, rollouts 6.
\end{itemize}

\subsection{Observability}
\texttt{raw\_simulation.csv} records the following fields:

\begin{lstlisting}[basicstyle=\ttfamily\scriptsize,frame=single]
run_id, world, strategy, seed, lifespan, alive, true_ruin,
passive_ruin, ruin_type, ruin_step, final_R, final_E, final_G,
final_O, final_K, final_X, cumulative_productivity,
peak_productivity, peak_exposure, mean_productivity
\end{lstlisting}

% =====================================================================
\section{Output Specification}
\label{sec:vnext-outputs}

\subsection{CSV Files}

\begin{table}[ht]
\centering
\small
\begin{tabularx}{\linewidth}{lX}
\toprule
\textbf{File} & \textbf{Content} \\
\midrule
\texttt{raw\_simulation.csv}
& Raw data of all episodes (per-episode record). \\
\texttt{world\_results.csv}
& Per-(world, strategy) aggregation + performance score. \\
\texttt{overall\_results.csv}
& Per-strategy cross-world aggregation. \\
\texttt{ruin\_analysis.csv}
& Detailed analysis of ruin causes. \\
\texttt{ruin\_summary.csv}
& Cross-world ruin-cause summary. \\
\bottomrule
\end{tabularx}
\caption[Output CSV files]{Output CSV files.}
\label{tab:vnext-csv-outputs}
\end{table}

\subsection{Visualisation}

\begin{table}[ht]
\centering
\small
\begin{tabularx}{\linewidth}{cXX}
\toprule
\textbf{\#} & \textbf{Plot} & \textbf{Type} \\
\midrule
1 & \texttt{performance\_comparison.png} & Composite score bar
chart. \\
2 & \texttt{survival\_by\_world.png} & Grouped bar (world $\times$
strategy). \\
3 & \texttt{exposure\_vs\_productivity.png} & Scatter plot. \\
4 & \texttt{optionality\_distribution.png} & Box plot. \\
5 & \texttt{ruin\_rates.png} & True vs passive ruin grouped bar. \\
6 & \texttt{lifespan\_distribution.png} & Box plot. \\
7 & \texttt{ruin\_cause\_breakdown.png} & Stacked bar (cause
composition). \\
8 & \texttt{survival\_heatmap.png} & Strategy $\times$ world
heatmap. \\
\bottomrule
\end{tabularx}
\caption[Visualisation outputs]{Visualisation outputs.}
\label{tab:vnext-plots}
\end{table}

% =====================================================================
\section{Module Structure}
\label{sec:vnext-modules}

\begin{table}[ht]
\centering
\small
\begin{tabularx}{\linewidth}{lX}
\toprule
\textbf{Module} & \textbf{Responsibility} \\
\midrule
\texttt{config/defaults.py}
& Centralised parameter management. \\
\texttt{worlds.py}
& World family definitions; per-run random draws. \\
\texttt{state.py}
& State vector; transition functions; ruin detection. \\
\texttt{nrmo\_core.py}
& NRMO veto logic (original + vNext); high-stakes mode. \\
\texttt{tuning\_layer.py}
& Adaptive tuning (manual / scenario / adaptive); hysteresis. \\
\texttt{strong\_engine.py}
& MC rollout search; candidate generation; scoring. \\
\texttt{strategies.py}
& Implementation of the 9 comparison strategies. \\
\texttt{simulator.py}
& Episode runner; full-factorial experiment. \\
\texttt{metrics.py}
& Aggregation; ruin analyser; composite score; CSV export. \\
\texttt{plots.py}
& Generation of 8 graph types. \\
\texttt{sweep.py}
& Parameter sweep. \\
\texttt{main.py}
& Entry point (smoke / full / sweep). \\
\bottomrule
\end{tabularx}
\caption[Module structure]{Module structure.}
\label{tab:vnext-modules}
\end{table}

% =====================================================================
\section{Execution Commands}
\label{sec:vnext-execution}

\begin{lstlisting}[basicstyle=\ttfamily\small,frame=single]
# Smoke test
python main.py --mode smoke

# Full experiment
python main.py --mode full

# Parameter sweep (specify world)
python main.py --mode sweep \
       --sweep-world Normal --sweep-samples 30

# Engine tuning
python main.py --mode smoke \
       --engine-rollouts 8 --engine-lookahead 5

# Custom output directory
python main.py --mode full --output-dir results_exp01
\end{lstlisting}

\paragraph{Dependencies.}
Python 3.10+, numpy, pandas, matplotlib.

% =====================================================================
\section{Future Extensions}
\label{sec:vnext-future-extensions}

\begin{enumerate}
\item Generation of complete narrative for the 30-civilisation
casebook.
\item Multi-agent interaction (inter-civilisation trade /
conflict / alliance).
\item Design of human-in-the-loop decision experiments.
\item Longitudinal operation-log studies (8--12-week real-world
decision logging).
\item Survival-time analysis via Cox proportional-hazard model.
\item Automation of parameter tuning via Bayesian optimisation.
\end{enumerate}

\medskip

\centering
\textit{--- End of Specification ---}\\
NRMO vNext + StrongEngine Civilisation Simulation v2.0\\
March 2026

% =====================================================================
\section*{Connections to Other Parts}

\begin{itemize}
\item This vNext specification is realised in the v2.5 engineering
implementation in Part~X.
\item The 6-dimensional civilisation state vector
(Section~\ref{sec:vnext-state}) and the 4-dimensional action vector
(Section~\ref{sec:vnext-action-space}) are the cognitive analogues
of the operational state-action structure used in the SOP layer
of the v4 monograph.
\item The strict separation principle
(Section~\ref{sec:vnext-separation}) implements the
governance--execution separation invariant
(Invariant~\ref{inv:gov-exec-sep}).
\item The 5 world families
(Section~\ref{sec:vnext-world-families}) feed the empirical
validation results in Part~XI.
\item The Adaptive Tuning Layer
(Section~\ref{sec:vnext-adaptive-tuning}) is extended in Part~IX
($\Omega$ Full v5.0) with the sustainable-growth estimator and
$\lambda_{\mathrm{drift}}$ penalty.
\end{itemize}
